\section{Methods}
\label{sec:methods}

\subsection{Evidence boundary and prospective discipline}

The evaluation separates three evidence classes. The complete-gold controlled
index tests mechanism behavior under an enumerable denominator. The protected
exact-contract battery isolates closure decisions when route state and
materiality are known. The TREC-COVID experiment tests the registered external
recall-and-cost claim on a public collection. Results from one class cannot
promote another: synthetic mechanism behavior does not establish external
retrieval superiority, and a secondary TREC-COVID metric cannot replace the
registered primary criteria.

For each result-bearing campaign, task identities, controller and comparator
identities, budgets, evaluator rules, decision thresholds, and repeat or
bootstrap seeds were frozen before outcomes were inspected. Candidate systems
received only public task views and label-free retrieved records. Complete gold
was available only to the evaluator. Post-outcome analysis could describe
failure stages or generate manuscript tables, but could not change the decision
rule.

The controller returns one of three scientific outcomes. A task passes when all
declared material obligations are discharged. It fails when the system makes a
checkable incorrect decision or exhausts the registered budget without a valid
closure. It remains undetermined when material evidence cannot be assessed,
such as when a required route is unavailable. An undetermined outcome is not a
zero and is not evidence of absence.

\subsection{Complete-gold controlled index}

The controlled index contains \OfflineTaskCount{} tasks derived from
\OfflineTopicCount{} topics and \OfflineDocumentCount{} synthetic documents.
Relevant items are distributed across lexical, semantic, citation,
reformulation, and restricted-provider routes. The generator includes shared
backends, duplicate locators, content revisions, extraction-question shifts,
and unavailable routes. Budgets are set below exhaustive probing so route
allocation and stopping remain observable.

Fourteen systems were frozen: the governed controller, six confirmatory
baselines, one exploratory adaptive comparator, and six ablations. Confirmatory
baselines include single-route lexical and dense retrieval, one-pass hybrid
retrieval, a protocol-driven systematic-review policy, and stopping controls.
All systems use the same host-owned session, route interface, resource meter,
and evaluator. Three deterministic repeat seeds test execution stability. They
do not increase the statistical unit beyond the \OfflineTaskCount{} tasks.

The main descriptive measures are complete-gold recall, premature task closure,
pass rate, undetermined rate, budget-exhaustion rate, duplicate processing,
legitimate rereading, route contribution, and route overlap. Route-stop false
positives and false negatives are evaluated separately from task closure. The
frozen precision plan requires 385 tasks for its committed tier; the assembled
\OfflineTaskCount{} tasks meet that count. However, the achieved half-width of 0.0496 exceeds
the frozen superiority margin of 0.03. The primary controlled comparison is
therefore labeled underpowered and cannot authorize a superiority claim.

\subsection{Exact-contract battery}

The protected battery contains 400 cases across five acquisition settings:
scientific literature, web or API evidence, dataset registries, code
repositories, and experiment or measurement routes. Cases distinguish complete
obligation discharge, one available route remaining open, material routes that
are unavailable, censored, or provider-invalid, duplicate coverage,
question-processing obligations, and low-utility routes that remain material.
Held-out variants add unavailable nonmaterial decoys, preventing a blanket rule
that any unavailable source blocks closure.

The primary comparator receives strong donor mechanisms: evidence coverage,
decision sufficiency, utility stopping, finite route certificates,
deduplication, route state, materiality, and question-conditioned processing.
It makes the same route-local decisions as the fail-closed controller. The only
systematic difference is task aggregation. After all available material routes
stop, the comparator excludes unavailable, censored, or provider-invalid
material routes from the closure denominator. A second comparator permits a
global sufficiency signal to close the task. A third comparator receives the
same unresolved-route semantics as the controller and defines an
information-equivalent boundary expected to tie.

The primary effect is the paired difference in exact scientific-closure
decisions between the fail-closed controller and the donor-complete available-
route product. The frozen practical margin is 0.10. Uncertainty uses a
domain-stratified bootstrap with 20,000 resamples, and discordance is reported
with exact McNemar counts. Nonregression requires zero false task closures,
perfect closure on clean cases, no unnecessary undetermined decisions, and no
materiality error on held-out decoys.

\subsection{Registered TREC-COVID comparison}

The external comparison uses the 50 TREC-COVID topics and the public test
collection \cite{voorhees2020treccovid}. Four arms were frozen: BM25, a
reciprocal-rank-fusion hybrid, the governed controller, and a diagnostic variant
that ignores unresolved route obligations when declaring completion. Each arm
is evaluated on the same topic set. The topic is the paired inference unit.

The registered pass rule is noncompensatory. First, recall@100 must be
noninferior to the best comparator with margin $-0.02$, assessed using the lower
bound of a 95\% paired bootstrap interval from 10,000 resamples. Second, the
controller must reduce candidates or tool calls by at least 25\%, with the
confidence interval's lower bound showing at least a 10\% reduction. Third,
premature-stop error must not increase. Failure of either recall or cost is
sufficient to fail the gate.

nDCG@10 was measured but was explicitly outside this pass rule. It describes
the ranking quality of the first ten results and can therefore move differently
from recall@100 and reading cost. We report it as secondary evidence only. The
registered gate is not amended after observing the favorable nDCG result.

\subsection{Bounded external stress tests}

Additional external probes are used for stage attribution and adverse-case
retention, not for the headline decision. A MetaSyn probe runs pinned BM25
retrieval and a deterministic public-protocol screening rule on all 86 released
test reviews, then applies the official identifier-only evaluator. A bounded
AutoResearchBench Deep probe uses the released title judge with a separate
nine-case positive and negative control. OpenAIRE/Crossref campaigns enforce a
provider-validity threshold before any comparison can be interpreted. Public
screening studies retain frozen active-learning donors and prospectively fixed
primary recall and work-saving criteria. Each probe keeps its own denominator,
authority, and failure state.
